
%%%%%%%%%%%%%%%%%%%%%%%%%%%%%%%%%%%
%%%%%%%%%%%%%%%%%%%%%%%%%%%%%%%%%%%
%%%%%%%%%%%%%%%%%%%%%%%%%%%%%%%%%%%
%%%%%%%%%%%%%%%%%%%%%%%%%%%%%%%%%%%

\newpage

\renewcommand{\arraystretch}{1.3}
\vspace*{-5em} 

\begin{center}
    \textbf{IIS di Esempio - Comune Esempio (XX)}\\
    \textbf{Griglia di Valutazione - Elaborato Scritto di Matematica - Liceo Scientifico}
\end{center}
\begin{tabularx}{\textwidth}{|m{6.8cm}|m{5.6cm}|m{3.5cm}|}
\hline
\textbf{DOCENTE:} Cognome Nome &
\multicolumn{2}{|l|}{\textbf{DATA DELLA VERIFICA:} \rule{1.15cm}{0.4pt}/\rule{1.15cm}{0.4pt}/\rule{1.25cm}{0.4pt}}\\
\hline
\textbf{COGNOME:} &
\textbf{NOME:} &
\textbf{CLASSE:} \\
\hline
\end{tabularx}
\\
\thispagestyle{empty}

\vspace{1mm}

\begin{tabularx}{\textwidth}{|>{\centering}m{3,5cm}|>{\centering\arraybackslash}m{11.6cm}|>{\centering\arraybackslash}m{0.8cm}|}
\hline
\textbf{INDICATORI} & \textbf{DESCRITTORI} & \textbf{PT}\\
\hline
\renewcommand{\arraystretch}{0.5}
&&\\
\end{tabularx}

\vspace{-0.75cm} % Spazio negativo per collegare le due tabelle

\renewcommand{\arraystretch}{1}
\begin{tabularx}{\textwidth}{|c|>{\arraybackslash}m{11.6cm}|>{\centering\arraybackslash}m{0.8cm}|}

\multirow{4}{*}{\parbox{3,5cm}{ \textbf{\\\\\\Comprendere} \vspace{0.5cm}\\ \fontsize{9}{2}\selectfont{Analizzare la situazione\\problematica, identificare\\i dati ed interpretarli. }}}
& \fontsize{8.5}{2}\selectfont{\textbf{Non comprende} le richieste o le recepisce in maniera \textbf{gravemente inesatta}, non 
riuscendo a individuare o a collegare i concetti chiave e le informazioni essenziali. } & 1\\
\cline{2-3}
& \fontsize{8.5}{2}\selectfont{Analizza ed interpreta le richieste in maniera \textbf{parziale}, riuscendo a selezionare \textbf{solo alcuni} dei concetti chiave e delle informazioni essenziali, o, pur avendone individuati diversi, \textbf{commette degli errori} nell'interpretarne alcuni e nello stabilire i collegamenti.} & 2 \\
\cline{2-3}
& \fontsize{8.5}{2}\selectfont{Analizza in modo \textbf{complessivamente adeguato} la situazione problematica, individuando e interpretando \textbf{correttamente} i concetti chiave, le informazioni e le relazioni tra queste.} & 3 \\
\cline{2-3}
& \fontsize{8.5}{2}\selectfont{Analizza in modo \textbf{adeguato} la situazione problematica, individuando e interpretando correttamente i concetti chiave, le informazioni e le relazioni tra queste; utilizza \textbf{quasi sempre con buona padronanza} i codici matematici grafico-simbolici.} & 4 \\
\cline{2-3}
& \fontsize{8.5}{2}\selectfont{Analizza ed interpreta in modo \textbf{completo e pertinente} i concetti chiave, le informazioni essenziali e le relazioni tra queste; utilizza i codici matematici grafico-simbolici con \textbf{buona padronanza e precisione}.} & \textbf{5} \\
\Xhline{2pt}

\multirow{4}{*}{\parbox{3,5cm}{ \textbf{\\\\\\Individuare} \vspace{0.5cm}\\ \fontsize{9}{2}\selectfont{Mettere in campo\\strategie risolutive e\\individuare la strategia\\più adatta.}}}
& \fontsize{8.5}{2}\selectfont{\textbf{Non individua} strategie di lavoro. \textbf{Non coglie alcuno spunto} nell'individuare il procedimento risolutivo.} & 0 \\
\cline{2-3}
& \fontsize{8.5}{2}\selectfont{Mette in campo strategie di lavoro \textbf{non adeguate}. \textbf{Non è in grado} di individuare relazioni corrette tra le variabili in gioco. \textbf{Non riesce} ad impostare correttamente le varie fasi del lavoro.} & 1 \\
\cline{2-3}
& \fontsize{8.5}{2}\selectfont{Attua strategie di lavoro \textbf{non sempre efficaci}, talora sviluppandole in modo \textbf{poco coerente}. Usa con una \textbf{certa difficoltà} le relazioni tra le variabili. Nella \textbf{maggior parte dei casi non riesce} ad impostare correttamente le varie fasi del lavoro.} & 2 \\
\cline{2-3}
& \fontsize{8.5}{2}\selectfont{Sa individuare delle strategie risolutive, anche se \textbf{non sempre le più adeguate ed efficienti}. Dimostra di conoscere le procedure consuete e le possibili relazioni tra le variabili e le utilizza in modo \textbf{adeguato}.} & 3 \\
\cline{2-3}
& \fontsize{8.5}{2}\selectfont{Attraverso congetture effettua \textbf{chiari collegamenti logici}. Individua strategie di lavoro \textbf{adeguate ed efficienti}. Utilizza in modo adeguato le relazioni matematiche note. Dimostra \textbf{padronanza} nell'impostare le varie fasi di lavoro.} & 4 \\
\cline{2-3}
& \fontsize{8.5}{2}\selectfont{Attraverso congetture effettua, \textbf{con padronanza}, chiari collegamenti logici. Individua strategie di lavoro \textbf{adeguate ed efficienti}. Dimostra padronanza nell'impostare le varie fasi di lavoro. Individua \textbf{con cura e precisione} le procedure \textbf{ottimali anche non standard}.} & \textbf{5} \\
\Xhline{2pt}

\multirow{4}{*}{\parbox{3,5cm}{ \textbf{\\\\\\Sviluppare il\\\\processo\\\\risolutivo}\vspace{0.5cm}\\ \fontsize{9}{2}\selectfont{Risolvere la situazione\\ problematica in maniera\\ coerente, completa\\e corretta, applicando\\le regole ed eseguendo\\i calcoli necessari.}}}
& \fontsize{8.5}{2}\selectfont{ \textbf{Non applica} le strategie scelte e \textbf{non sviluppa} il processo risolutivo. \textbf{Non riesce} ad applicare procedure e/o teoremi. \textbf{Non ottiene alcuna soluzione}.} & 1 \\
\cline{2-3}
& \fontsize{8.5}{2}\selectfont{Applica le strategie scelte, \textbf{per lo più, in maniera non corretta}: sviluppa \textbf{la maggior parte} 
del processo risolutivo in modo \textbf{incompleto e/o errato}. Applica nella \textbf{maggior parte dei casi} procedure e/o teoremi in modo \textbf{errato e/o con diversi errori} nei calcoli. } & 2 \\
\cline{2-3}
& \fontsize{8.5}{2}\selectfont{ \textbf{Quasi sempre} applica le strategie scelte in maniera corretta pur presentando delle \textbf{imprecisioni} e sviluppando il processo risolutivo \textbf{quasi completamente}. È in grado di utilizzare procedure e/o teoremi e di applicarli \textbf{quasi sempre} in modo corretto e appropriato. Commette \textbf{alcuni errori} nei calcoli.} & 3 \\
\cline{2-3}
& \fontsize{8.5}{2}\selectfont{ \textbf{Quasi sempre} applica le strategie scelte in maniera corretta e sviluppa il processo risolutivo in modo \textbf{completo}. Applica \textbf{quasi sempre} procedure e/o teoremi o regole in modo corretto e appropriato. Esegue i calcoli in modo \textbf{quasi corretto}.} & 4 \\
\cline{2-3}
& \fontsize{8.5}{2}\selectfont{ Applica le strategie scelte in maniera \textbf{corretta}, sviluppa il processo risolutivo in modo \textbf{corretto, chiaro e completo}. Applica procedure e/o teoremi o regole in modo corretto e appropriato, \textbf{con abilità}. Esegue i calcoli in modo \textbf{accurato}.} & \textbf{5} \\
\Xhline{2pt}

\multirow{4}{*}{\parbox{3,5cm}{ \textbf{\\\\\\Argomentare}\vspace{0.5cm}\\ \fontsize{9}{2}\selectfont{\mbox{Commentare  e giustificare}\\opportunamente la scelta\\della strategia applicata,\\i passaggi fondamentali\\del processo esecutivo e\\la coerenza dei risultati. }}}
& \fontsize{8.5}{2}\selectfont{\textbf{Non argomenta} la strategia/procedura risolutiva.} & 1 \\
\cline{2-3}
& \fontsize{8.5}{2}\selectfont{\textbf{Non argomenta} o, per lo più, argomenta in modo \textbf{errato} la strategia/procedura risolutiva utilizzando un linguaggio matematico \textbf{non appropriato o molto impreciso}.} & 2 \\
\cline{2-3}
& \fontsize{8.5}{2}\selectfont{Argomenta in maniera \textbf{frammentaria e/o non sempre coerente} la strategia/procedura esecutiva. Utilizza un linguaggio matematico \textbf{per lo più appropriato}, ma \textbf{non sempre rigoroso}.} & 3 \\
\cline{2-3}
& \fontsize{8.5}{2}\selectfont{Argomenta in modo \textbf{coerente ma incompleto} la procedura esecutiva. Utilizza un linguaggio matematico \textbf{pertinente ma con qualche incertezza}.} & 4 \\
\cline{2-3}
& \fontsize{8.5}{2}\selectfont{Argomenta in modo \textbf{coerente, preciso ed esaustivo} le strategie adottate. Mostra una \textbf{buona padronanza} nell'utilizzo del linguaggio scientifico.} & \textbf{5} \\
\Xhline{2pt}
\end{tabularx}

%[End-griglia]

\vspace{0.2cm}
\hspace{14.2cm} TOTALE: \hspace{1cm}/20
\vspace{0.5cm}

Comune Esempio,\raisebox{-1mm}{\rule{1cm}{0.4pt}}/\raisebox{-1mm}{\rule{1cm}{0.4pt}}/2026
\hspace{0.5cm} Docente:\rule{5cm}{0.4pt}
\hspace{2.1cm} Voto: \hspace{1.7cm}/10